\clearpage
\SourceLeaf{66}
\pdfbookmark[1]{Disquisitio de elementis ellipticis Palladis — Anzeige 1810}{anzeige-palladis1810}
\begin{center}\rule{.75\textwidth}{.3pt}\\
Göttingische gelehrte Anzeigen. 1810. December 13.\\
\rule{.75\textwidth}{.3pt}\end{center}
Am 25. November übergab Herr Prof. \textsc{Gauss} der königl. Societät der Wissenschaften
eine Vorlesung:
\begin{center}
\textit{Disquisitio de elementis ellipticis Palladis ex oppositionibus annorum}\\
1803, 1804, 1805, 1807, 1808, 1809.
\end{center}
Seit drei Jahren hatte der Verfasser keine neue Rechnungen über die Bahn
der Pallas angestellt. Die Beobachtungen vom Jahre 1808 waren bei der grossen
Lichtschwäche des Planeten sehr dürftig und mangelhaft gewesen, und sind zum
Theil erst später bekannt geworden, daher es nicht der Mühe werth schien, schon
damals die Elemente darnach zu verbessern. Erst nachdem Herr Prof. \textsc{Gauss} die
Beobachtungen der Pallas von 1809, welche von \textsc{Bouvard} auf der kaiserl. Sternwarte
in Paris angestellt waren, erhalten hatte, berechnete er nebst der Opposition
von 1809 zugleich die von 1808 (s. diese Gel. Anz. 1810 Febr. 24). Die fernere
Discussion aller bisher beobachteten sechs Oppositionen ergab das Resultat,
dass eine elliptische Bahn nicht mehr zureicht, sie alle genau darzustellen: eine
Folge der grossen Störungen, die dieser Planet von den übrigen, und besonders
von dem Jupiter, erleidet. Herr Prof. \textsc{Gauss} hat dies auf eine doppelte Art gezeigt.
Zuvörderst berechnete er drei Systeme von elliptischen Elementen, jedes
aus vier Oppositionen, nemlich das erste aus denen von 1803, 1804, 1805, 1807;
\SourceLeaf{67}%
das zweite aus denen von 1804, 1805, 1807, 1808; das dritte aus den Oppositionen
von 1805, 1807, 1808, 1809, unter denen sich nur kleine Verschiedenheiten
hätten zeigen müssen, wenn die Bewegung rein elliptisch wäre. Wir stellen
hier diese drei Systeme neben einander:
\begin{center}
\setlength{\tabcolsep}{5pt}
\begin{tabular}{c|c|c|c}
& I. & II. & III.\\\hline
$A$ & $221^0\;39'\;30''4$ & $221^0\;34'\;56''7$ & $221^0\;23'\;24''6$\\
$B$ & $770.2143$ & $770.4467$ & $770.9265$\\
$C$ & $121\quad3\quad11.4$ & $121\quad5\quad22.1$ & $120\;58\quad4.8$\\
$D$ & $172\;28\;56.9$ & $172\;28\;46.8$ & $172\;27\;52.4$\\
$E$ & $34\;37\;41.0$ & $34\;37\;31.5$ & $34\;36\;49.4$\\
$F$ & $0.2450198$ & $0.2447624$ & $0.2446335$\\
$G$ & $0.4423149$ & $0.4422276$ & $0.4420473$
\end{tabular}
\end{center}
wo
\begin{description}\setlength{\itemsep}{0pt}\setlength{\parsep}{0pt}\setlength{\parskip}{0pt}
\item[$A$] Epoche der mittlern Länge 1803 für den Meridian von Göttingen
\item[$B$] mittlere tägliche tropische Bewegung
\item[$C$] Länge der Sonnennähe 1803
\item[$D$] Länge des aufsteigenden Knotens 1803
\item[$E$] Neigung der Bahn
\item[$F$] Excentricität
\item[$G$] Logarithm der halben grossen Axe
\end{description}
bedeutet. So nothwendig es nun sein wird, die Störungen der Pallas durch die
andern Planeten mit in Rechnung zu bringen, wenn man eine bleibende Übereinstimmung
der Erscheinungen mit der Rechnung beabsichtigt, so wird man doch
auch künftig zu manchen Zwecken eine rein-elliptische Bahn vorziehen, wenn
man dadurch den unerträglich weitläuftigen Rechnungen ausweichen kann, die
die grosse Anzahl der Störungsgleichungen sonst nothwendig machen würde. Namentlich
ist es, wenn es bloss die Wiederauffindung des Planeten gilt, gewiss
bequemer, rein-elliptische Elemente auf die zunächst vorhergegangenen vier Oppositionen
zu gründen, und darnach die Ephemeride zu construiren, welche nach
einem Jahre noch nicht sehr viel vom Himmel abweichen kann. Deshalb hat
Herr Prof. \textsc{Gauss} auch die im Octoberheft der Mon. Corresp. abgedruckte Ephemeride
für die Bewegung der Pallas in ihrer nächsten Erscheinung nach dem
obigen III. System von Elementen berechnet.

\SourceLeaf{68}%
Die zweite Art, wie Herr Prof. \textsc{Gauss} den Einfluss der Störungen nachgewiesen
hat, besteht in der Berechnung von rein-elliptischen Elementen, die sich
an \textit{alle} sechs Oppositionen \textit{möglichst} genau anschliessen, und die dessen ungeachtet
sich von den einzelnen beobachteten Örtern bedeutend entfernen. Wir setzen
auch dieses \textit{vierte} System von Elementen hier her:
\begin{center}
\begin{tabular}{cll}
$A$ & $.\;.\;.\;.\;.\;.\;.\;.\;.\;.\;.\;.\;.\;.\;.\;.$ & $221^0\;34'\;53''64$\\
$B$ & $.\;.\;.\;.\;.\;.\;.\;.\;.\;.\;.\;.\;.\;.\;.\;.$ & $770.5010$\\
$C$ & $.\;.\;.\;.\;.\;.\;.\;.\;.\;.\;.\;.\;.\;.\;.\;.$ & $121\quad8\quad8.54$\\
$D$ & $.\;.\;.\;.\;.\;.\;.\;.\;.\;.\;.\;.\;.\;.\;.\;.$ & $172\;28\;12.43$\\
$E$ & $.\;.\;.\;.\;.\;.\;.\;.\;.\;.\;.\;.\;.\;.\;.\;.$ & $34\;37\;28.35$\\
$F$ & $.\;.\;.\;.\;.\;.\;.\;.\;.\;.\;.\;.\;.\;.\;.\;.$ & $0.2447424$\\
$G$ & $.\;.\;.\;.\;.\;.\;.\;.\;.\;.\;.\;.\;.\;.\;.\;.$ & $0.4422071$
\end{tabular}
\end{center}
Die Fehler dieser Elemente stellt folgende Übersicht dar:
\begin{center}\small
\setlength{\tabcolsep}{4pt}
\begin{tabular}{c|r|r}
Opposition & \multicolumn{2}{c}{Unterschied}\\
von & der heliocentr. Länge & der geocentr. Breite\\\hline
1803 & $-111''00$ & $-8''31$\\
1804 & $+59.18$ & $-36.67$\\
1805 & $+19.92$ & $+0.07$\\
1807 & $+85.77$ & $+25.01$\\
1808 & $+135.88$ & $+28.72$\\
1809 & $-216.54$ & $+83.01$
\end{tabular}\end{center}
Den Hauptinhalt der vorliegenden Abhandlung macht die Entwickelung der
Methoden aus, wie diese verschiedenen Systeme von Elementen gefunden sind,
und in dieser Rücksicht kann dieselbe als eine Art von Supplement zu einigen
Abschnitten der \textit{Theoria motus corporum coelestium} betrachtet werden. Die Aufgabe,
aus vier beobachteten Örtern eines Himmelskörpers seine Bahn zu bestimmen,
war in jenem Werke zwar schon umständlich abgehandelt, und zwei Auflösungen
gegeben, eine für den Fall, wo die Bahn noch ganz unbekannt ist, und
eine für den Fall, wo die schon näherungsweise bekannte Bahn verbessert werden
soll. Letztere hätte also auch zur Bestimmung obiger drei Systeme von Elementen
in Anwendung gebracht werden können; allein gerade im vorliegenden
Fall, wo die vier Beobachtungen Oppositionen sind, hat Herr Prof. \textsc{Gauss} es vortheilhafter
gefunden, eine ganz andere Methode zu gebrauchen, welche in der
\SourceLeaf{69}%
Abhandlung ausführlich erklärt, und durch das Beispiel der wirklichen Berechnung
des dritten Systems der Elemente noch mehr erläutert ist. Eine nähere
Beschreibung dieser Methode hier zu geben verstattet der Raum nicht.

Die Berechnung des vierten Systems von Elementen ist nach den Grundsätzen
geführt, die in dem 3. Abschnitt des 2. Buchs der \textit{Theoria motus corporum
coelestium} entwickelt sind, und die vorliegende Abhandlung gibt auch hiezu mehrere
Zusätze, die hoffentlich den Astronomen nicht unwillkommen sein werden.
Zuerst eine bequeme Berechnung der Differential-Änderungen der heliocentrischen
Länge und der geocentrischen Breite aus den Differential-Änderungen der
einzelnen Elemente. Sodann ein eigenes Verfahren, die unbekannten Grössen
dem oben erwähnten Grundsatze gemäss zu bestimmen. Sind nemlich $w,w',w''$
u. s. w. die vorgegebenen linearen Functionen der unbekannten Grössen $p,q,r$
u. s. w., und soll das Aggregat $ww+w'w'+w''w''+$ u. s. w. ein Kleinstes werden,
so erhält man leicht so viele lineäre Gleichungen, als unbekannte Grössen
sind, aus denen diese durch Elimination bestimmt werden müssen. Diese Elimination
ist aber, wenn die Anzahl der unbekannten Grössen etwas beträchtlich
ist, eine äusserst beschwerliche Arbeit, und zwar deswegen, weil jede der Gleichungen
alle unbekannten Grössen enthält. Herr Prof. \textsc{Gauss} hat diese Arbeit
sehr bedeutend abgekürzt; denn obgleich er die Auflösung auch auf so viele lineäre
Gleichungen, als unbekannte Grössen sind, zurückführt, so sind diese
Gleichungen so beschaffen, dass nur die erste alle unbekannten Grössen enthält,
aber die zweite von $p$, die dritte von $p$ und $q$; die vierte von $p,q$ und $r$ frei
ist u. s. w., daher die Bestimmung der unbekannten Grössen in der umgekehrten
Ordnung nur noch wenige Mühe macht. Ausserdem hat diese Methode noch den
Vortheil, dass man den kleinsten Werth von $ww+w'w'+w''w''+$ u. s. w. im voraus
angeben, und so die Vergleichung desselben mit dem nachher berechneten,
wenn in $w,w',w''$ etc. die für die unbekannten Grössen gefundenen Werthe substituirt
werden, zu einer Controlle der Rechnung benutzt werden kann.
\par\bigskip\centerline{\rule{.24\textwidth}{.3pt}}
